
\chapter{Resource-aware long-term memory on a smartphone}
\label{ch:on_device_memory}

On the tested Galaxy S23 Ultra, Vector RAG had higher accuracy point estimates and lower measured
energy per query than truncated full context in all five reader pairs. These separate measurements
describe answer quality and device cost. Correct answers per kilojoule summarizes their joint
effect: the observed Vector-RAG advantage was 5.8$\times$ to 26.9$\times$ for the four reader pairs
with a non-zero full-context denominator, and these multipliers are uncertain with 20 questions
per cell (Section~\ref{sec:ch5_charge}). On this slice, Vector RAG answered every question that full
context answered correctly, as well as additional questions. Across the five readers, it produced
40.8 to 68.8 correct answers per full-charge equivalent, against 0 to 8.9 for full context. Query
service was also 6.8$\times$ to 8.2$\times$ faster. These values describe the tested phone
workloads, not every mobile agent or question mix.

\paragraph{Contribution of this chapter.} Among the studies reviewed in
Section~\ref{sec:ch2_device}, MobileRAG reports per-query phone energy for document retrieval and
SmartRAG reports average phone power for one graph system whose baselines ran off the phone. None
compares memory representations of a conversation history by accuracy and energy both measured on
one handset. This chapter makes that comparison for vector retrieval and truncated full context
with five readers, using accuracy judged from the phone's own answers. It also measures, on the same
handset and with the same energy method, what it costs to build a vector memory and an
SLM-extracted graph memory for a conversation.

The result continues the strategy comparison in Chapter~\ref{ch:memory_evaluation} but
changes the unit of analysis. Chapter~\ref{ch:memory_evaluation} asks which memory
representation produces correct answers under controlled workstation conditions. This
chapter asks what happens when the same broad choice must fit a phone's battery, memory
and thermal limits. Chapter~\ref{ch:spe} later returns to graph-memory repair. Here,
GraphRAG is limited to the direct phone construction evidence that actually exists.

The chapter addresses the following research question:

\begin{quote}
\textbf{RQ2.} What trade-offs between answer quality and resource consumption arise when
deploying long-term memory strategies on smartphones?
\end{quote}

\section{A resource-aware objective}
\label{sec:ch5_objective}

A continuing mobile agent does more than issue isolated language-model calls. It stores
observations, updates its memory and retrieves evidence while sharing a finite battery
and thermal envelope with the rest of the phone. A memory design that is accurate but
repeatedly fills the context window can be too slow or expensive to use. Low energy is
not sufficient either if the reader returns no correct answers.

Accuracy and energy per query are reported first as separate measures of answer quality and device
cost. Correct answers per unit of energy, $G_E$ in Equation~\ref{eq:golden_metric}, adds a combined
measure: how many correct answers the measured energy buys. It is reported per kilojoule as the
number of judge-scored correct phone answers divided by the voltage-weighted battery energy of the
recorded workload, including energy spent on incorrect answers. The combined measure complements
the separate results; a high value alone does not establish adequate accuracy or low query cost.
Correct answers per full-charge equivalent, $G_Q$, are computed separately from the direct
charge-counter delta.

This avoids mixing a voltage-weighted joule numerator with a nominal-voltage battery
denominator. Neither measure predicts daily battery life because the screen, radios,
background applications and the rest of the agent loop sit outside the measured span.

Section~\ref{sec:ch2_device} explains why these costs must be measured for the tested
configuration rather than inferred from token counts or from other
devices~\citep{slm_edge_deployment,ondevice_lm_review}.

\section{Experimental boundary}
\label{sec:ch5_boundary}

Chapter~\ref{ch:methodology} defines the phone protocol, voltage-weighted Coulomb
counting, question selection and judging procedure. Inference ran on the CPU of the Galaxy
S23 Ultra described there, with 12~GB of RAM and a 5,000~mAh battery. The local inference server used one slot, an
8,192-token slot context and five model threads. Workloads ran unplugged with the screen
off.

Five standard Q4\_K\_M readers were evaluated: Qwen3.5 0.8B, 2B and 4B, and Gemma 4 E2B
and E4B. Vector retrieval used the Q4\_K\_M Nomic Embed Text v1.5 artifact. Revisions, filenames, byte sizes and
SHA-256 digests reconstructed during the evidence audit are
recorded in \texttt{docs/runbooks/MODEL\_PROVENANCE.md}. The surviving run directories
did not embed immutable model hashes, so this reconstruction combines server logs with
the surviving model cache.

The primary comparison uses the same balanced 20-question LoCoMo slice under Vector RAG
and truncated full context~\citep{locomo}. It contains five questions from each of the
four evaluated question families. Decoding was deterministic, and every accuracy in the
phone tables was judged from a phone output. The energy unit is a complete sitting
divided by its number of scored answers. Vector RAG has three complete energy sittings
for every reader. Gemma E2B full context also has three, one of which completed 19
answers. The other full-context cells have one complete sitting. Appendix~\ref{app:on_device_statistics} gives the replicate values and uncertainty details.

The device audit searched the experiment locations accessible through ADB and Termux:
shared storage, Termux's private results directory and the local mirror of pulled
results. Protected Android sandboxes and restricted kernel interfaces were not readable
without root. The evidence claim is therefore limited to accessible experimental
storage.

\section{The two statistical results used in this chapter}
\label{sec:ch5_statistical_story}

\subsection{Workstation accuracy cannot substitute for phone accuracy}
\label{sec:ch5_parity}

The phrase ``the workstation result does not transfer to the phone'' has a precise
meaning in this thesis. It does not mean that the phone was proved less accurate. It
means that the predefined agreement rule did not license the two environments as
interchangeable.

The parity run used the same 100 question identifiers and the same E2B Vector-RAG
configuration on the workstation and phone. Each answer was reduced to a paired binary
correctness label. The two devices agreed on 83 of 100 labels, below the required 90\%
agreement. Cohen's $\kappa$ was 0.621. Of the 17 disagreements, five questions were
correct only on the workstation and 12 only on the phone. An exact McNemar test on those
discordant pairs gave $p=0.1435$.

These numbers answer different questions. Agreement measures whether the same questions
receive the same labels. McNemar's test asks whether the disagreements consistently
favor one device; this run did not establish a directional accuracy shift. The agreement
criterion still failed. Workstation accuracy is therefore not joined to phone energy anywhere in
the main result.

\subsection{Question family changes the observed difficulty}
\label{sec:ch5_question_families}

The second result concerns the benchmark, not the device. The agreement rule
forbids substituting workstation \emph{accuracy} for phone accuracy; it does not prevent
using the workstation record to characterize the \emph{questions} both environments
answer. That use assumes only that the relative difficulty of LoCoMo's question families
carries across environments, which is weaker than the per-question label agreement the
parity gate tested and failed. On the workstation sample, question family (the LoCoMo
category) remained associated with correctness in all eight primary Gemma reader-memory
cells after Holm correction (Section~\ref{sec:ch4_locomo_categories}). Phone accuracy
must therefore be tied to its sampled question composition. The balanced phone slice
gives each family equal weight and should not be read as an estimate under LoCoMo's
natural family frequencies.

\section{Measured phone comparison}
\label{sec:ch5_phone_comparison}

Table~\ref{tab:ch5_phone_cells} presents the direct phone measurements.\footnote{All
20-question phone Vector RAG runs in this chapter, including
Figures~\ref{fig:ch5_phone_efficiency} and~\ref{fig:ch5_resource_dashboard}, use the
unprefixed embedding path. On the same 20 questions, the prefixed 100-question runs answered 8
(E2B) and 10 (E4B) correctly, against 10 and 8 here (Section~\ref{sec:retrieval}).} Repeated cells
use the mean across complete sittings. Single-sitting values are point estimates without
a replication claim.

\begin{table}[htbp]
  \centering
  \small
  \setlength{\tabcolsep}{4.2pt}
  \begin{tabular}{llrrrr}
    \toprule
    Reader & Memory & Accuracy & J/answer & Query time (s) & Sittings \\
    \midrule
    Qwen3.5 0.8B & Vector RAG   & 20\% & 306  & 15.0  & 3 \\
    Qwen3.5 0.8B & Full context &  0\% & 330  & 114.5 & 1 \\
    Qwen3.5 2B   & Vector RAG   & 30\% & 433  & 34.8  & 3 \\
    Qwen3.5 2B   & Full context & 10\% & 850  & 262.0 & 1 \\
    Qwen3.5 4B   & Vector RAG   & 35\% & 569  & 94.0  & 3 \\
    Qwen3.5 4B   & Full context & 10\% & 2166 & 641.9 & 1 \\
    Gemma 4 E2B  & Vector RAG   & 50\% & 551  & 54.5  & 3 \\
    Gemma 4 E2B  & Full context &  5\% & 1395 & 420.7 & 3$^{a}$ \\
    Gemma 4 E4B  & Vector RAG   & 40\% & 737  & 93.0  & 3 \\
    Gemma 4 E4B  & Full context &  5\% & 2589 & 759.7 & 1 \\
    \bottomrule
  \end{tabular}
  \caption{Phone-measured accuracy, voltage-weighted energy and query latency on the
  balanced 20-question LoCoMo slice. Energy is the mean joules per answer across complete
  sittings where repeated measurements exist. Latency is the mean of the per-sitting
  medians. $^{a}$One Gemma E2B full-context sitting completed 19 answers.}
  \label{tab:ch5_phone_cells}
\end{table}

Vector RAG had a higher accuracy point estimate, lower energy per answer and shorter
query time in every reader pair. For Qwen3.5 0.8B the energy difference is small, 306 against
330~J, and lies within the spread of that reader's three vector sittings (248 to 345~J); for the
other four readers full context used 2 to 4 times more energy per answer. Reader size alone did not order the systems. Qwen3.5
0.8B used the least energy for a Vector-RAG answer but reached only 20\% accuracy. Gemma
E2B reached the highest measured accuracy and correct answers per unit of energy, while the larger Gemma
E4B used more energy and scored ten percentage points lower on this slice. The small
sample does not support a general ranking of the readers. It does show why a phone
policy cannot use parameter count as its only decision variable.

\begin{figure}[htbp]
  \centering
  \includegraphics[width=\textwidth]{results/thesis_visuals/fig05_phone_efficiency.pdf}
  \caption{Correct answers per unit of energy and energy per answer across the five phone
  readers. The left panel divides
  phone-measured accuracy by replicate-mean energy; the right panel shows mean joules per
  answer, with each cell's individual sittings as grey dots above its marker. The Qwen3.5 0.8B full-context
  cell produced zero correct answers, so no finite Vector/full-context efficiency ratio
  is assigned to that pair.}
  \label{fig:ch5_phone_efficiency}
\end{figure}

\section{Correct answers per charge}
\label{sec:ch5_charge}

Table~\ref{tab:ch5_charge} converts the Vector-RAG charge-counter measurements into
correct answers per 5,000 mAh full-charge equivalent. It also reports the observed
within-sitting advantage over full context.

\begin{table}[htbp]
  \centering
  \small
  \begin{tabular}{lrr}
    \toprule
    Reader & Vector correct/full charge & Paired Vector/FC efficiency \\
    \midrule
    Qwen3.5 0.8B & 50.2 & Undefined$^{a}$ \\
    Qwen3.5 2B   & 52.8 & 5.8$\times$ \\
    Qwen3.5 4B   & 46.7 & 13.4$\times$ \\
    Gemma 4 E2B  & 68.8 & 25.1$\times$ \\
    Gemma 4 E4B  & 40.8 & 26.9$\times$ \\
    \bottomrule
  \end{tabular}
  \caption{Correct answers per full-charge equivalent for Vector RAG and its observed advantage over
  full context. Vector values use mean charge per answer across complete sittings. The
  efficiency ratios are computed within matched sittings and are therefore not obtained
  by dividing the rounded means in Table~\ref{tab:ch5_phone_cells}. $^{a}$Full context
  produced zero correct answers, so the ratio is undefined.}
  \label{tab:ch5_charge}
\end{table}

The range in the opening paragraph is now visible directly in Table~\ref{tab:ch5_charge}: Gemma E4B is the 40.8 floor and Gemma E2B is the 68.8 ceiling.
The smallest reader does not maximize correct answers per charge because its lower accuracy offsets
part of its energy saving. The largest Gemma reader also does not maximize it because
its extra energy did not buy a higher score on this slice.

The complete full-context sittings produced 8.9 correct answers per charge for Qwen3.5
2B, 3.3 for Qwen3.5 4B, 2.4 in the primary Gemma E2B sitting, 1.4 for Gemma E4B and zero
for Qwen3.5 0.8B. These are workload-normalized comparisons, not a recommendation to
drain a handset or a forecast of daily use.

\paragraph{Sensitivity of the efficiency ratio to small-sample denominators.}
The $25.1\times$ and $26.9\times$ ratios rest on very small denominators. On a 20-question slice, a
5.0\% full-context accuracy is one correct answer ($1/20$), with a Wilson 95\% confidence interval of
$[0.9\%, 23.6\%]$. One more correct full-context answer ($2/20=10.0\%$) would have cut the E2B ratio from
$25.1\times$ to about $12.5\times$, and zero correct answers, as observed for Qwen 0.8B, leave the ratio
undefined. Resampling the 20 observed question pairs with mean energy held fixed gives
2.5th-percentile ratios ranging from 2.6$\times$ to 8.8$\times$ across the four readers with defined
ratios (Appendix~\ref{app:ch5_energy_replication}). These are conditional sensitivity results:
they do not include between-sitting energy variation or charge-counter bias, and resampling cannot
introduce full-context-only successes absent from the original sample. The observed efficiency
ordering holds in all five reader pairs, but the multipliers describe this slice and should not
be read as general constants.

\section{Latency, RAM and thermal headroom}
\label{sec:ch5_resources}

Figure~\ref{fig:ch5_resource_dashboard} places accuracy and energy beside query latency,
the peak temperature of the hottest thermal zone and peak system RAM.

\begin{figure}[htbp]
  \centering
  \includegraphics[width=\textwidth]{results/thesis_visuals/fig12_phone_workload_dashboard.pdf}
  \caption{Phone workload comparison across five readers. Accuracy uses 20 balanced LoCoMo questions with
  Wilson 95\% intervals. Resource whiskers show between-sitting standard deviations; grey dots
  are individual sittings. Vector retrieval has three sittings per reader; full context has
  three for E2B and one otherwise. Energy and latency axes are logarithmic; thermal and RAM axes
  do not start at zero. Resources summarize complete run spans. Temperature is the maximum
  across readable thermal zones, not battery temperature; RAM is system-wide. Neither peak
  isolates a model component.}
  \label{fig:ch5_resource_dashboard}
\end{figure}

Vector retrieval shortened the mean of the per-sitting median query times from 114.5 to
15.0 seconds for Qwen3.5 0.8B, from 262.0 to 34.8 seconds for Qwen3.5 2B, from 641.9 to
94.0 seconds for Qwen3.5 4B, from 420.7 to 54.5 seconds for Gemma E2B and from 759.7 to
93.0 seconds for Gemma E4B. Retrieval adds ingestion, but repeated queries avoid paying
full-context prefill on every request.

Peak RAM did not follow the same ordering. Measured as system-wide memory in use (total minus
available), the cell means lay between about 5.5 and 6.7 GB, with
retrieval lower for some readers and higher for others. Model weights and the resident
runtime occupy much of this footprint at the tested context size. The evidence supports
reduced repeated context processing, not a claim that retrieval always reduces peak
resident memory.

The displayed mean of per-sitting hottest-zone peaks was higher for Vector RAG in all five reader
pairs, by 0.5 to 3.7 degrees Celsius, despite much shorter query times. Four comparisons
combine a three-sitting Vector-RAG mean with one full-context sitting, so this 5-of-5
pattern is descriptive. It does not establish that retrieval inherently runs hotter, and the recorded traces do not
explain it. Vector-RAG queries were short, with medians of 15 to 94 seconds, against about 2 to
13 minutes for full context, so a peak reading may reflect a short burst rather than cumulative
thermal load. However, prefill and decode were not timed separately, governor decisions were not
recorded, and the core-frequency medians discussed below show no consistent pattern between the
strategies. Testing the burst account would require a controlled experiment with matched starting
temperature and frequency-residency logging.

Mean power over a sitting separates the two strategies as well, and in the opposite direction to
duration. Table~\ref{tab:ch5_power} gives the whole-device average over each measured span. Vector
retrieval drew more average power than full context in four of the five reader pairs, by 0.13 to
1.22~W, and the same power to within 0.01~W in the fifth, while its median query took about one
seventh as long.

\begin{table}[htbp]
  \centering
  \small
  \begin{tabular}{lrr}
    \toprule
    Reader & Vector RAG (W) & Full context (W) \\
    \midrule
    Qwen3.5 0.8B & 2.88 & 2.89 \\
    Qwen3.5 2B   & 3.97 & 3.27 \\
    Qwen3.5 4B   & 3.44 & 3.31 \\
    Gemma 4 E2B  & 4.11 & 3.59 \\
    Gemma 4 E4B  & 4.71 & 3.49 \\
    \bottomrule
  \end{tabular}
  \caption{Mean whole-device power over each sitting's measured span, from the same charge-counter
  energy and span as the energy column of Table~\ref{tab:ch5_phone_cells}
  (\texttt{results/golden\_metric.json}). One sitting per cell, not a replicate mean. The
  construction pilot of Section~\ref{sec:ch5_graph_boundary} averaged 3.96~W over its span.}
  \label{tab:ch5_power}
\end{table}

Two explanations were considered for that difference (\texttt{results/phone\_power\_phases.json}).
The first is the work of building the vector memory. Each vector sitting embeds six conversation
histories, one before each history's questions, and these embedding windows took 22.8 to 31.7 minutes
of a 29 to 55-minute sitting. They drew 3.4 to 5.1~W, more than the question stretches of the same
sitting in four of the five readers (by 0.4 to 0.9~W; Qwen3.5 4B was level), so building the memory
accounts for part of the difference. The question stretches alone drew 3.1 to 4.4~W, against 2.9 to
3.6~W over the whole full-context sittings: a gap of 0.9~W for Gemma E4B and 0.3~W or less for the
other readers. Four of those full-context sittings started at the highest possible charge (100\%),
and in three of them the charge counter paused, which makes their power read low, so only the E4B
gap is clearly larger than that effect. The second explanation is the duty cycle. The median of the maximum core frequency
during the question stretches was 1,594 to 1,978~MHz in the vector sittings, against 1,843~MHz in
every full-context sitting, so the clock data show no consistent pattern that would support it. The
traces do not record governor decisions, so the remaining difference is not explained here. Vector
stretches that began while the charge counter was paused at full charge are left out of these phase
figures.

The recovered second Gemma E4B full-context sitting illustrates another device outcome:
failure to complete. It produced 8 of 20 output records and no final summary during
approximately 113 minutes. The battery level fell from 96\% to 65\% over the surviving
trace. Because the boundary cannot be matched to the sealed cells, this run is excluded
from the result tables. Its value is operational: checkpointing and recoverability are
part of a phone memory system, not merely conveniences for the experimenter.

\subsection{Prefill dominance and KV-cache size}
\label{sec:ch5_prefill_decode_memory}

The latency and memory measurements in the dashboard reflect the prefill--decode difference and the
growth of the KV cache with context length, both introduced in Section~\ref{sec:ch2_slms}.
Figure~\ref{fig:ch5_prefill_decode_memory} shows the measured query times of the Gemma~4 E2B phone runs
(llama.cpp on the Snapdragon 8 Gen 2 CPU cluster, 5 pinned threads, CPU mask \texttt{0xf8}).

\begin{figure}[htbp]
  \centering
  \includegraphics[width=\textwidth]{results/thesis_visuals/fig20_phone_prefill_decode_memory.pdf}
  \caption{Measured query time (retrieval plus generation) of every answered Gemma-4 E2B question on the
  Samsung Galaxy S23 Ultra against its prompt length: three full-context sittings, four vector-retrieval
  sittings with ten retrieved turns (one of them 100 questions long), one sitting with two retrieved turns, and
  the four questions of the early graph pilot. Large markers are medians. The phone records only the total
  generation time of each question, so prefill and decode time are not separated.}
  \label{fig:ch5_prefill_decode_memory}
\end{figure}

Figure~\ref{fig:ch5_prefill_decode_memory} plots every answered E2B phone question against its prompt length.
Full-context prompts of about 4,900 tokens took a median of 411 to 431 seconds per sitting. Ten retrieved
turns (about 680 tokens) took 46 to 65 seconds, two retrieved turns (about 290 tokens) about 21 seconds, and
the four graph-pilot questions (about 510 tokens) about 34 seconds. Query time follows prompt length: the
median full-context prompt is 17 times as long as the median two-turn prompt (4,905 against 294 tokens) and
took 20 times as long (418.5 against 20.8 seconds). The answers are short, a median of 10 tokens (8 to 21)
across the 243 answers. The phone logs only the total generation time, so the time spent decoding the answer
cannot be separated in these runs; the near-proportional scaling indicates that processing the prompt sets
the cost of a question. Measured per sitting, the phone processed 10.9 to 15.4 prompt tokens per second of
query time, and 11.0 to 11.9 in the full-context sittings. Averaged over sittings, full context took 7.7
times as long per question as vector retrieval, 420.7 against 54.5 seconds (Section~\ref{sec:ch5_resources}).

Memory was measured as system-wide RAM in use (total minus available memory), not per process. Peak use in
the E2B sittings averaged about 6.7~GB under both strategies (Figure~\ref{fig:ch5_resource_dashboard}). The
server reserves the key--value cache for its whole 8,192-token context when it loads, so the reservation was
the same for every strategy: prompt length changed time, not memory. The size of that cache follows from the
model files. Twenty of Gemma-4 E2B's 35 layers reuse the KV cache of earlier layers, and 12 of the 15 layers
that keep their own cache use a 512-token sliding window, so only three global layers grow with context.
Calculated from the GGUF header shapes, the E2B cache is 57~MB at 8,192 tokens and 0.81~GB at 131,072 tokens,
against 4.4~GB still free at the highest measured E2B peak (Figure~\ref{fig:app_ch5_kv_wall}). The layout
behind that calculation is the published one: the Gemma~4 technical report states a four-to-one ratio of
local to global attention blocks for E2B and five-to-one for E4B, and KV-cache sharing across 20 of 35 layers
in E2B and 18 of 42 in E4B~\citep{gemma4}, which are the counts used here. Its Table~3 also reports the
footprint of an int8 cache at a 32,768-token context, 0.05~GB for E2B and 0.14~GB for E4B, about half of what
the calculation above gives for the same contexts at half precision. The calculated figures are therefore
treated as an upper bound on the cache; the comparison against free memory only becomes more favourable if
the published footprints are the operative ones. Apart from a four-question pilot that doubled the
history budget to 8,192 tokens and was not pursued further (Appendix~\ref{app:reproducibility}), no
longer context was run on the phone and no long-context study was done on it; the long-window
full-context runs of Section~\ref{sec:ch4_long_fc} were made on the workstation. The runtime's other
buffers were not measured at longer contexts, so this comparison bounds the cache alone rather than
showing that a longer context would fit.
Bounded KV growth does not make a long context cheap either, because prefill time and energy still grow with
every prompt token, as the query times show.

\section{The measured graph-construction boundary}
\label{sec:ch5_graph_boundary}

Graph memory remains preliminary on the phone. Figure~\ref{fig:ch5_phone_graph_pilot} shows the one
direct construction pilot. Building an SLM-extracted graph for 419 dialogue turns took 92.1 minutes
and consumed 21.9 kJ under the same voltage-weighted integration used elsewhere in this chapter, or
52.3 J per turn. Fifty-eight of 84 extraction batches failed. The pilot then answered zero of four
questions correctly. Four answers do not support a comparative accuracy estimate.

The timeline in Figure~\ref{fig:ch5_phone_graph_pilot} shows how the handset behaved during
extended extraction. The temperature trace is the hottest of all readable thermal zones. Its first
sample, taken after the model servers had loaded, read $47.9^\circ\text{C}$, with the battery at
$33.3^\circ\text{C}$. The hottest zone peaked at $95.4^\circ\text{C}$ 74 seconds after launch, 29 seconds
into extraction, and the Cortex-X3 prime core reached its 3.36~GHz maximum only within the first minute.
For most of the 92.1-minute construction span the fastest core then ran at 1.59~GHz, and the hottest zone
had a median of $74.6^\circ\text{C}$, with ten-minute medians between $73.4$ and $77.8^\circ\text{C}$ after
the first ten minutes. The charge counter gives an average of 3.96~W over the span; with the clock held
down, the heat input and with it the temperature stayed nearly constant. A steady clock ceiling under a
CPU-bound workload is consistent with thermal or power throttling, but the trace does not record governor
state, so it cannot confirm the cause. Battery temperature rose steadily from $33.3^\circ\text{C}$ to
$44.3^\circ\text{C}$. The two traces diverge throughout, which is why hottest-zone peaks are not reported
as battery or chassis temperatures.

The construction cost is the most informative measurement from this pilot.
Table~\ref{tab:ch5_construction} shows how it is compared with building a vector memory. Each side
is the whole-device battery energy spent while the phone built memory, measured with the same
charge-counter method as the rest of this chapter, divided by the number of dialogue turns built.
Graph construction is the Gemma~4 E2B reader writing triples for batches of five turns. Vector
ingestion is the Nomic embedder embedding each turn once. Neither side includes answering questions.

\begin{table}[htbp]
  \centering
  \small
  \begin{tabular}{p{0.27\linewidth}p{0.30\linewidth}p{0.33\linewidth}}
    \toprule
    & Graph construction & Vector ingestion \\
    \midrule
    Work done per turn & Triple extraction by Gemma~4 E2B, five turns per call & One embedding by Nomic
    Embed v1.5 \\
    Measured runs & One pilot build & Three Gemma E2B vector sittings \\
    Histories built & One conversation (\texttt{conv-26}, 419 turns) & Six per sitting; two windows
    left out$^{a}$ \\
    Duration and mean power & 92.1 min at 3.96~W & 31.7 min at 4.29~W in the first sitting \\
    Energy per dialogue turn & 52.3~J & 2.04~J (1.85 to 2.37 by sitting) \\
    \midrule
    Ratio & \multicolumn{2}{l}{About 26$\times$ (22$\times$ to 28$\times$ against each vector sitting)} \\
    \bottomrule
  \end{tabular}
  \caption{How the construction ratio is computed. The graph column is the single early phone build
  (256-token extraction limit, 58 of 84 batches failed), not a usable repaired graph, and the ratio
  applies to that build only. Both columns are direct phone measurements over the build windows only (\texttt{results/golden\_metric.json}, \texttt{construction} entry;
  duration and power from \texttt{results/phone\_power\_phases.json}). $^{a}$Two embedding windows
  began while the charge counter was paused at full charge and are left out.}
  \label{tab:ch5_construction}
\end{table}

The two builds drew a similar average power, so the ratio comes almost entirely from time. Writing
triples took about 13 seconds per dialogue turn (92.1 minutes for 419 turns), while embedding took
about half a second per turn in the first vector sitting (31.7 minutes for 3,435 turns). Neither
side depends on judged accuracy, on token counts or on the workstation record.

The ratio describes this one attempted build, not the cost of constructing a usable graph, and three
limits apply to it. The graph side is a single build of one conversation, so it has no replicate
spread; the 22--28$\times$ range reflects only the variation between vector sittings. The build used
the early 256-token extraction limit, and 58 of its 84 batches failed; failed batches still spent
their decoding energy, so they are included in the cost, but they contributed no triples. Finally,
the repaired 512-token configuration of Chapter~\ref{ch:graph_mechanisms} was not built on the phone,
so its construction energy was not measured.

The asymmetry matters because a memory policy that inspects only query prompts cannot observe
either term: construction is paid before the first question is asked.

\begin{figure}[htbp]
  \centering
  \includegraphics[width=\textwidth]{results/thesis_visuals/fig06_phone_graphrag_pilot.pdf}
  \caption{Early graph-construction pilot on the Galaxy S23 Ultra: 419 turns, 92.1 minutes and 21.9~kJ
  (52.3~J/turn), with 58 of 84 extraction batches failing. The red trace is the hottest readable
  thermal zone, not battery temperature; sampling begins after server loading. Core frequencies
  and charge use describe the recorded workload. The four subsequent answers do not constitute a
  comparative graph-QA evaluation.}
  \label{fig:ch5_phone_graph_pilot}
\end{figure}

The pilot supports one narrow conclusion: a memory policy must account for construction
as well as query service, as in the accounting of Equation~\ref{eq:ch2_amortization}. It
does not support a phone routing comparison among graph,
vector and full-context memory because no complete phone GraphRAG QA cell exists.
Figure~\ref{fig:ch5_graph_build_scaling} places the build time of this directly measured 419-turn phone
build ($21.9\,\text{kJ}$, $52.3\,\text{J/turn}$) beside workstation build times across dialogue lengths,
which grow linearly with the number of turns. The figure shows time only; the energy values are the
ones given here.

\begin{figure}[htbp]
  \centering
  \includegraphics[width=\textwidth]{results/thesis_visuals/fig11_graph_build_scaling.pdf}
  \caption{Graph construction time. Left: 50 workstation builds across ten LoCoMo histories and five
  Gemma extractor configurations (E2B v1--v3; E4B v1 and v3), with a linear fit per
  configuration. Right: fitted workstation times at 419 turns and the directly measured
  92.1-minute phone build, on a logarithmic axis. The early phone pilot differs from the
  repaired workstation configurations.}
  \label{fig:ch5_graph_build_scaling}
\end{figure}

The earlier token-routing grid is therefore removed from the main argument. Appendix~\ref{app:on_device_statistics} records why it is exploratory and keeps the long-context
and amortization calculations separate from direct measurements.

\section{Position relative to mobile systems literature}
\label{sec:ch5_literature}

Section~\ref{sec:ch2_device} reviews the mobile studies closest to this chapter and the gap
they leave. Two of them report measurements that bear directly on the results above.
SmartRAG reports average power of 2.37 to 4.77~W during triplet extraction and 2.12 to
5.36~W during question answering on its two OnePlus handsets, but not energy per
answer~\citep{smartrag}. It also finds that retrieval-conditioned prefill dominates latency:
on a Snapdragon 8 Elite Gen 5 with the 1.7B backbone, median time to first token was 1.32~s
for question decomposition but 36.9~s for answering, consistent with the prompt-driven query
times observed here in Section~\ref{sec:ch5_prefill_decode_memory}. MobileRAG reports real-phone energy
for document RAG on a Galaxy S24, from 24.71~J per query to more than 200~J in some
configurations~\citep{mobilerag_energy}.
In both cases the handset, model, corpus, retrieval work and measurement span differ, so
their numbers are not pooled with the measurements in this chapter.

\section{Implications for resource-aware long-term memory}
\label{sec:ch5_implications}

The measurements support four design choices for a phone-native agent.

\begin{enumerate}
  \item \textbf{Optimize correct answers per unit of energy rather than one counter.} Energy, latency and
  accuracy changed together. Correct answers per kilojoule and per charge reveal cases
  in which the lowest-energy reader is not the most useful one.

  \item \textbf{Use bounded retrieval as the default path for the tested workload.}
  Vector RAG occupied a better measured accuracy, energy and latency point than full
  context in all five reader pairs. The conclusion is limited to this retriever, device
  and balanced question slice.

  \item \textbf{Keep a lifecycle ledger.} Embedding, graph extraction, prefill, decode
  and persistent storage are different operations. Prompt tokens describe reader input
  but omit memory construction entirely, and on this phone construction is where the
  strategies separate most sharply: the early graph build cost about 26 times more energy per
  turn than embedding (Table~\ref{tab:ch5_construction}). A ledger that prices only the query
  path misses that difference before the first question is served.

  \item \textbf{Treat completion and provenance as system properties.} A run that
  consumes battery but cannot be resumed has delivered no sealed work. Checkpoints,
  model digests, runtime identity and explicit measurement spans are necessary if a
  memory policy is to learn from device telemetry.
\end{enumerate}

On a phone, ``long term'' does not mean only a larger store or a longer prompt. It means
recovering useful history across repeated interactions without exhausting the resources
that allow the agent to continue operating.

\section{Experimental limitations}
\label{sec:ch5_threats}

\paragraph{One device and runtime.}
The measurements come from one Galaxy S23 Ultra, one Android and kernel build and a
CPU-oriented llama.cpp configuration. A newer SoC, accelerator support, another thread
count or another operating-system governor may change the ordering.

\paragraph{Small phone sample and uneven replication.}
Each accuracy cell contains 20 balanced questions. Question family is associated with
correctness in the larger workstation sample, but the phone slice is too small for
family-specific phone claims. Vector energy has three sittings for every reader, while
full-context replication is complete only for Gemma E2B. Appendix~\ref{app:on_device_statistics} reports this asymmetry instead of hiding it in main-text
error bars.

\paragraph{Short-workload overhead.}
In the 20-question Gemma E2B vector sittings, building the memory for six histories took about
60\% of the energy (Appendix~\ref{app:ch5_run_length}). A separate 100-question E2B Vector-RAG
run, which spread ten histories over 100 questions, used less energy per answer, but it used the
task prefixes and its question composition and execution time differ. It is evidence that the
number of questions per history matters, not a matched estimate of a universal amortization factor.

\paragraph{Automated judging.}
Phone answers are scored by the same automatic judge as the workstation answers, with the
limits stated in Section~\ref{sec:judge}.

\paragraph{Model and binary provenance.}
Exact model revisions and digests were reconstructed from the surviving cache and
server logs. The per-run directories did not preserve immutable hashes. The installed
server identifies itself as \texttt{0.1.2-dev}, build zero, with an unknown commit. A
source checkout survives, but it cannot be asserted to be the binary's build commit.

\paragraph{Incomplete and inaccessible evidence.}
The incomplete E4B sitting is excluded from the sealed tables. The audit covered all
experimental locations accessible through ADB and Termux, but not protected Android
sandboxes or root-only kernel interfaces.

\paragraph{Graph evidence.}
Only one early phone graph construction pilot was measured, and no complete phone graph
QA cell exists. The main chapter makes no graph-routing or graph-accuracy claim.

\section{Chapter conclusion}
\label{sec:ch5_conclusion}

The phone results support a direct systems conclusion. Vector retrieval produced higher
accuracy point estimates, lower energy per answer and much shorter query times than
full context under all five readers. It delivered 40.8 to 68.8 correct answers per
full-charge equivalent, and its finite observed advantage over full context ranged from
5.8$\times$ to 26.9$\times$ correct answers per kilojoule. This combined measure summarizes the
separate quality and energy results.

Two statistical results set the limits of that claim. The workstation-to-phone parity
gate failed, so phone accuracy must be measured on the phone. Question family was also
associated with correctness in every primary Gemma cell, so the balanced phone result
must remain attached to its question mix. RAM and thermal peaks do not reduce to the
same ordering as energy and latency, and the graph pilot shows that construction can be
a large part of the memory lifecycle. A resource-aware memory policy therefore needs these
device measurements before it can route between strategies.
